% Fragmentos leídos directamente de analysis/analizar.py mediante listings.
\section{Implementación del script de análisis}
\label{sec:implementacion-analisis}

El archivo \texttt{analysis/analizar.py} transforma los CSV de latencia
en resúmenes, tablas y figuras. Recibe los archivos de cada mecanismo
por separado mediante \texttt{--sockets} y \texttt{--grpc}. Cada fila
representa un mensaje completado. El formato de entrada sigue pendiente
de contrastarse con los CSV reales de las implementaciones.

\subsection{Lectura y validación de los registros}
La función \texttt{leer\_csv} utiliza pandas para leer los campos
inicialmente como texto. La codificación \texttt{utf-8-sig} permite
leer también un CSV que tenga una marca BOM:

\lstinputlisting[firstline=41,lastline=41]{../analysis/analizar.py}

El lector también adapta el formato nativo de seis columnas de sockets.
Se indican el modo con \texttt{--sockets} o
\texttt{--sockets-concurrente} y, si está documentado, el equipo con
\texttt{--equipo-sockets}. Una duración de corrida ausente en ese formato
deja el rendimiento como N/D.
La validación rechaza campos vacíos en los datos requeridos, modos distintos de secuencial o
concurrente y etiquetas de calentamiento diferentes de 0 y 1.
El tamaño del mensaje y la latencia deben ser positivos y finitos.
La duración debe cumplir ese criterio en el formato completo; además,
el tamaño debe ser entero. El mecanismo se añade a los
registros a partir del argumento usado al cargar el archivo.

En \texttt{resumir} se revisa que cada identificador sea único dentro de
su corrida y que todas sus filas tengan la misma duración y etiqueta
de calentamiento. Si se encuentra una inconsistencia, el programa
termina con un mensaje de error.

\subsection{Calentamiento y estadísticas de latencia}
Las filas marcadas con 1 corresponden al calentamiento. Para calcular
las estadísticas se seleccionan las que tienen la etiqueta 0:

\lstinputlisting[firstline=101,lastline=101]{../analysis/analizar.py}

La agrupación \texttt{GRUPO} contiene mecanismo, equipo, modo y tamaño
en bytes. Así se conserva una fila de resultados por condición.
La latencia se expresa en milisegundos y se resume mediante:

\lstinputlisting[firstline=105,lastline=111]{../analysis/analizar.py}

La media utiliza todos los mensajes válidos del grupo; la mediana
representa su posición central. \texttt{std} calcula la desviación
estándar muestral, con divisor $n-1$. El percentil 95 se obtiene con
la interpolación lineal predeterminada de \texttt{quantile}.
Si solo existe una observación, la desviación muestral queda sin definir
y se presenta como N/D.

\subsection{Rendimiento y duración de las corridas}
La clave \texttt{CORRIDA} añade el identificador de repetición a las
variables de agrupación. Como la duración total está repetida en cada
fila del CSV, se toma una sola vez antes de sumar los tiempos:

\lstinputlisting[firstline=112,lastline=120]{../analysis/analizar.py}

El rendimiento agregado divide los mensajes completados entre ese
tiempo total. En concurrencia se necesita la duración común de los tres
clientes, desde el comienzo de la carga hasta que termine el último.
La suma de las latencias individuales mide otra cantidad, porque
las solicitudes pueden solaparse.

\subsection{Tablas y gráficos generados}
La opción \texttt{--informe} crea una carpeta nueva con
\texttt{resumen.csv}, \texttt{por\_corrida.csv}, dos tablas LaTeX y
los gráficos de cajas y barras en PNG y PDF.
\texttt{escribir\_tabla} convierte los resultados a fragmentos LaTeX
y escapa caracteres especiales de los identificadores.

El diagrama de cajas usa las latencias individuales. Para las barras,
\texttt{resumen\_corridas} calcula primero la media de cada corrida.
El siguiente fragmento dibuja su media y la desviación estándar entre
corridas:

\lstinputlisting[firstline=227,lastline=230]{../analysis/analizar.py}

Las dos desviaciones tienen alcances distintos: la tabla describe la
variación entre mensajes y la barra de error describe la variación
entre las medias de las corridas. Cada barra indica cuántas corridas
se utilizaron. Si solo hay una, se omite la barra de error.

\subsection{Estado de la comprobación}
El script se comprobó inicialmente con datos sintéticos temporales.
El 19 de noviembre se aplicó a los cuatro CSV reales recibidos y se
contrastaron sus estadísticas con NumPy. Se conservaron las filas
originales y se verificó que sus hashes SHA-256 permanecieran iguales
después del análisis. Los resultados se presentan en la
sección~\ref{sec:resultados}.
